% ============================================================
% 语义回响：通过回收被丢弃Token嵌入增强语言模型表达能力
% ============================================================

\documentclass[a4paper, 11pt]{article}
\usepackage[UTF8]{ctex}
\usepackage[utf8]{inputenc}
\usepackage[T1]{fontenc}
\usepackage[top=2.5cm, bottom=2.5cm, left=2.8cm, right=2.8cm]{geometry}
\usepackage{setspace}
\usepackage{amsmath, amssymb, amsthm}
\usepackage{bm}
\usepackage{graphicx}
\usepackage{booktabs}
\usepackage{tabularx}
\usepackage{array}
\usepackage{caption}
\usepackage{subcaption}
\usepackage[sort&compress, numbers]{natbib}
\usepackage{hyperref}
\hypersetup{colorlinks=true,citecolor=blue,linkcolor=black,urlcolor=blue}
\usepackage{microtype}
\usepackage{url}
\usepackage{threeparttable}
\usepackage{multirow}
\usepackage{enumitem}
\usepackage{algorithm2e}
\usepackage{xcolor}
\RestyleAlgo{plain}\SetAlgoCaptionSeparator{~}\SetAlgoVlined
\graphicspath{{图片/}}
\newcommand{\methodname}{Semantic Echo}
\newcommand{\lambdasym}{\lambda}
\newcommand{\gammasym}{\gamma}
\newtheorem{definition}{定义}
\newtheorem{theorem}{定理}
\renewcommand{\baselinestretch}{1.15}

\begin{document}

\title{\vspace{-1.2cm}
  \LARGE \textbf{语义回响：通过回收被丢弃Token嵌入\\增强语言模型表达能力}\\
  \large \textbf{Semantic Echo: Enhancing LLM Expressiveness by Recycling Discarded Token Embeddings}
  \vspace{-0.3cm}}
\author{
  邓斯键${}^{1\dagger}$（思路）\qquad
  DeepSeek${}^{2\ddagger}$（AI实现）\\[4pt]
  \small${}^\dagger$核心概念、技术路线与实验设计 \\
  \small${}^\ddagger$代码实现、实验执行与论文撰写}
\date{\small 2026年7月}
\maketitle
\thispagestyle{empty}

\begin{abstract}
\noindent 自回归语言模型在每一步解码中仅选取单个Token，其余候选Token的隐藏状态被完全丢弃本文称此为\textbf{概率暴政}（Probabilistic Tyranny）。
受人类交流中"潜台词"与"回响"现象的启发，本文提出\textbf{语义回响}（Semantic Echo）范式：
回收被丢弃Token的隐藏状态，经随机静态投影后重新注入当前表示空间，使模型感知更丰富的语义上下文。
我们进一步提出了\textbf{情感词库筛选}与\textbf{思考阶段分离注入}两项创新机制。
在Qwen2.5-0.5B-Instruct上的三轮共13组对照实验表明：
（1）$\lambda=0.5$弱回响条件下，原始方案可使语义熵提升+17.9\%；
（2）情感筛选机制可有效缓解$\lambda\geq1.0$时的重复生成问题（E8 vs E4: +25.33\%）；
（3）三种保留策略对比表明，衰减策略最优，滑动窗口其次，全局保留最弱；
（4）语义回响不修改模型权重、不重新训练，仅通过采样层扩展实现。

\vspace{0.3cm}
\noindent \textbf{关键词：} 大型语言模型，语义回响，Token嵌入回收，情感筛选，思考阶段分离，可控文本生成
\end{abstract}
